\documentclass[11pt,a4paper]{article}
\usepackage[T1]{fontenc}
\usepackage[utf8]{inputenc}
\usepackage{lmodern}
\usepackage{amsmath,amssymb,amsthm}
\usepackage{geometry}
\usepackage{microtype}
\usepackage{hyperref}
\usepackage{enumitem}
\geometry{left=23mm,right=23mm,top=23mm,bottom=23mm}
\hypersetup{
  hidelinks,
  pdftitle={Continuous Planar Isotropy and the Pythagorean Formula: A Lean-Formalized Explanation},
  pdfauthor={ChatGPT (GPT-6, High thinking effort); prompted by Alessandro Odasso},
  pdfsubject={AI-generated explanatory manuscript with explicit research-prompt disclosure}
}
\newtheorem{theorem}{Theorem}[section]
\newtheorem{lemma}[theorem]{Lemma}
\newtheorem{remark}[theorem]{Remark}
\setlist[itemize]{itemsep=2pt,topsep=4pt}
\setlist[enumerate]{itemsep=4pt,topsep=5pt}
\title{\textbf{Continuous Planar Isotropy and the Pythagorean Formula}\\[0.35em]
\large A Plain-English Explanation of Lean-Formalized Hypotheses and Proofs}
\author{AI-generated mathematical exposition: \textbf{ChatGPT (GPT-6, High thinking effort)}\\
Research question and initiating prompt: \textbf{Alessandro Odasso}}
\date{}
\begin{document}
\maketitle

\begin{abstract}
Why should the length of a displacement with coordinates $(x,y)$ be
$\sqrt{x^2+y^2}$? In this manuscript, length begins as an arbitrary
function, not as that familiar formula. We explain two related
implications formalized in Lean~4: a conditional Pythagorean theorem
using rigid motions, reflection, and axis reachability; and a stronger
characterization in which continuity of a linear planar action
\emph{derives} axis reachability and reflection symmetry. An explicit
trigonometric model proves the converse, while the maximum norm shows
why finite rotational symmetries alone are insufficient. The proofs
are conditional statements about explicitly assumed linear motions;
they do not construct geometry from incidence axioms or establish
historical novelty.
\end{abstract}

\section{Research prompt, authorship, and scope}

\noindent\textbf{AI and prompt disclosure.}
Alessandro Odasso initiated the research question by asking OpenAI's
ChatGPT whether it could find an \emph{original proof of the
Pythagorean theorem}. The manuscript text and proposed mathematical
argument were generated with ChatGPT (GPT-6, High thinking effort);
the accompanying formalization and subsequent extensions were
developed with AI assistance and checked in Lean~4/Mathlib. The
word \emph{original} reports Odasso's objective, not a verified
claim of priority, novelty, or originality in the mathematical
literature. Lean verifies formal implications from stated hypotheses,
not such historical claims.

This text is an accessible explanation of what the repository
actually formalizes. Its logical starting point is not a triangle
diagram or the distance formula: it is an unknown function of two
coordinates and an explicitly described family of linear
transformations. The key mathematical point is to keep \emph{length}
separate from the coordinate expression that the proof eventually
forces.

\section{The objects and assumptions}

Represent a displacement by $(x,y)\in\mathbb R^2$. Let
\[
 L:\mathbb R^2\longrightarrow\mathbb R
\]
be a \emph{primitive length function}. Initially there is no
definition $L(x,y)=\sqrt{x^2+y^2}$, no norm axiom, and no
assumed continuity of $L$.

Each planar map below is linear and is represented by a real
$2\times 2$ matrix. Write $F_t$ for a one-parameter family of
such maps and $Q$ for a distinguished quarter-turn. The assumptions
of the main characterization are:
\begin{enumerate}
 \item \textbf{Action law.} For all $s,t$, $F_{s+t}=F_s\circ F_t$.
 \item \textbf{Primitive oriented quarter-turn.}
       $Q(1,0)=(0,1)$ and $Q^2(x,y)=(-x,-y)$ for all $x,y$.
       The matrix of $Q$ is \emph{not} assumed in advance.
 \item \textbf{Membership.} There is a real parameter $\theta$
       for which $F_\theta=Q$.
 \item \textbf{Continuous action coefficients.}
       The first-column entries
       $t\mapsto (F_t)_{11}$ and
       $t\mapsto (F_t)_{21}$ are continuous.
 \item \textbf{Axial calibration.}
       $L(z,0)=|z|$ for every real $z$.
 \item \textbf{Rigidity.}
       $L(F_t(x,y))=L(x,y)$ for every $t,x,y$.
\end{enumerate}
These premises assume a sufficiently rich \emph{linear} motion
action; they do not infer its existence from incidence or
congruence alone. In contrast, we do \emph{not} assume that
$L$ is positive, continuous, homogeneous, subadditive, or
symmetric under reflection. The latter properties will follow
from identifying $L$.

\begin{theorem}[Characterization formalized in Lean]
\label{thm:main}
For every function $L$ satisfying axial calibration,
\[
\begin{split}
 &\text{there exist }F,Q\text{ satisfying the action,
 quarter-turn,}\\[-1pt]
 &\text{membership, continuity, and rigidity hypotheses}\\
 &\qquad\Longleftrightarrow\qquad
 L(x,y)=\sqrt{x^2+y^2}\quad\text{for all }x,y.
\end{split}
\]
In particular, under those hypotheses every displacement with
$c=L(a,b)$ satisfies $a^2+b^2=c^2$.
\end{theorem}

The equivalence is the Lean theorem
\texttt{continuous\_isotropy\_iff\_euclidean\_length}.
We explain its forward and reverse implications in the order
used by the proof modules.

\section{Forward implication: five proof steps}

\subsection*{Step 1. The quarter-turn and the shape of every motion}

Applying the square law to $(1,0)$ gives
\[
 Q(0,1)=Q^2(1,0)=(-1,0).
\]
Linearity then forces
\[
 Q(x,y)=(-y,x),\qquad
 J=\begin{pmatrix}0&-1\\1&0\end{pmatrix}.
\]
This matrix is a \emph{conclusion}, not a premise.
By the action law and $F_\theta=Q$,
\[
 F_tQ=F_{t+\theta}=F_{\theta+t}=QF_t.
\]
A real $2\times2$ matrix commuting with $J$ must have
the form
\[
 F_t=\begin{pmatrix}u(t)&-v(t)\\v(t)&u(t)\end{pmatrix},
 \qquad
 F_t(x,y)=(u(t)x-v(t)y,\;v(t)x+u(t)y).
 \tag{1}\label{eq:normal}
\]
The action law and the invertibility of $Q$ also give $F_0=I$,
so $u(0)=1$, $v(0)=0$.
These deductions correspond to
\texttt{primitive\_quarter\_turn\_is\_J},
\texttt{primitive\_quarter\_turn\_forces\_flow\_shape}, and
\texttt{action\_zero\_of\_primitive\_quarter\_turn}.

\subsection*{Step 2. Continuity excludes radial scaling}

Define the nonnegative quantity
\[
 s(t)=u(t)^2+v(t)^2.
\]
Multiplying matrices of the form~\eqref{eq:normal} yields
\[
 s(t+r)=s(t)s(r),\quad s(0)=1,\quad s(\theta)=1.
 \tag{2}
\]
Consequently $s(t+\theta)=s(t)$. The parameter $\theta$
cannot be zero because $F_0=I$ whereas $Q(1,0)=(0,1)$.
By the coefficient continuity, $s$ is continuous with a
nonzero period, hence has a compact range and a greatest
value $M\ge 1$.

Choose $r$ with $s(r)=M$. For every $t$,
\[
 s(t)M=s(t+r)\le M,
\]
so $s(t)\le1$. But (2) also gives
$s(t)s(-t)=s(0)=1$, and applying the same bound
to $-t$ implies $s(t)\ge1$. Therefore
\[
 \boxed{u(t)^2+v(t)^2=1\quad\text{for every }t.}
 \tag{3}
\]
Expanding the transformed coordinates now shows
\[
 (ux-vy)^2+(vx+uy)^2
 =(u^2+v^2)(x^2+y^2)=x^2+y^2. \tag{4}
\]
This does \emph{not} use $L$, its reflection symmetry, or any
derivative. In Lean the main ingredients are
\texttt{periodic\_nonnegative\_character\_eq\_one} and
\texttt{continuous\_action\_preserves\_quadratic}.

\subsection*{Step 3. Every displacement reaches the positive axis}

For a fixed displacement $(a,b)$, consider the second coordinate
of $F_t(a,b)$. By~\eqref{eq:normal} it is the continuous function
\[
 f(t)=v(t)a+u(t)b.
\]
At $t=0$ its value is $b$. At $t=2\theta$,
the action gives $F_{2\theta}=Q^2=-I$, so its
value is $-b$. The intermediate value theorem therefore
provides a parameter $T$ with $f(T)=0$.
Write $F_T(a,b)=(z,0)$. If $z<0$, composing with
the half-turn $F_{2\theta}$ changes $(z,0)$
to $(-z,0)$. Thus in every case
\[
 \exists T,z\quad z\ge0,\qquad F_T(a,b)=(z,0).
 \tag{5}
\]
The formal result is
\texttt{axis\_reachable\_of\_continuous\_action}.
This is a derived fact, not an extra axiom in the
continuous theorem.

\subsection*{Step 4. Compare the two descriptions of the axis point}

Equation~(4) applied to~(5) gives
\[
 z^2=a^2+b^2.
\]
On the other hand, rigidity and axial calibration say
\[
 L(a,b)=L(F_T(a,b))=L(z,0)=|z|=z.
\]
Because $z\ge0$,
\[
 \boxed{L(a,b)=\sqrt{a^2+b^2}},\qquad
 \boxed{a^2+b^2=L(a,b)^2}. \tag{6}
\]
The length formula has been \emph{proved} under the
action assumptions, rather than used to define $L$.
This is \texttt{length\_eq\_sqrt\_without\_mirror}.
The identity in~(6) is the coordinate
Pythagorean formula for the displacement whose primitive
length is $c=L(a,b)$.

\subsection*{Step 5. Mirror symmetry is a consequence}

Formula~(6) immediately yields
\[
 L(x,-y)=\sqrt{x^2+(-y)^2}
        =\sqrt{x^2+y^2}=L(x,y).
\]
Thus mirror symmetry, unlike in the earlier conditional proof,
is not a hypothesis. Lean records this as
\texttt{mirror\_symmetry\_of\_continuous\_isotropy}.

\section{Reverse implication: an explicit model}

The reverse direction does not merely assert that such
motions ought to exist. For
$L(x,y)=\sqrt{x^2+y^2}$, set
\[
 F_t=\begin{pmatrix}
 \cos t&-\sin t\\
 \sin t& \cos t
 \end{pmatrix},\qquad Q=F_{\pi/2}.
\]
The trigonometric addition identities prove the action law;
sine and cosine are continuous; $Q$ has the required
quarter-turn behavior; and the usual identity
$\sin^2t+\cos^2t=1$ proves rigidity.
The Lean definitions \texttt{standardRotation},
\texttt{standardQuarterTurn}, and the theorem
\path{continuous_isotropy_model_of_euclidean_length}
formalize this converse. This construction uses ordinary real
analysis; it is not a derivation of trigonometry from synthetic
geometry.

\section{How the earlier conditional theorem fits}

The archived Lean theorem
\texttt{pythagoras\_from\_metric\_free\_quarter\_turn}
does \emph{not} require continuity of the action.
Instead, in addition to the action, quarter-turn,
membership, axial calibration, and rigidity,
it explicitly assumes:
\begin{itemize}
 \item \emph{Mirror symmetry:} $L(x,-y)=L(x,y)$;
 \item \emph{Axis reachability for the selected vector:}
       $F_T(a,b)=(z,0)$ for some $z\ge0$;
 \item \emph{Length identification:} $c=L(a,b)$.
\end{itemize}
Its proof obtains the normal form~\eqref{eq:normal}
from the quarter-turn. Instead of Step~2's compactness
argument, mirror symmetry and rigidity show
\[
 L(u,v)=1,\quad L(u,-v)=1,\quad
 L(u^2+v^2,0)=1.
\]
Calibration and $u^2+v^2\ge0$ imply
$u^2+v^2=1$. With the \emph{assumed} axis
reachability, the algebraic calculation in Step~4 yields
$a^2+b^2=c^2$.

The two theorems therefore trade assumptions:
the older theorem permits a discontinuous action but
requires mirror symmetry and an axis hit; the newer
theorem requires continuity of two action coefficients
and \emph{derives} both consequences. This is not a claim
that every single assumption in one statement is weaker
than the corresponding assumption in the other.

\section{Why one quarter-turn is not enough}

Consider the maximum norm
\[
 B(x,y)=\max\{|x|,|y|\}.
\]
It is continuous, axially calibrated, positive definite,
homogeneous, and satisfies the triangle inequality.
It is also invariant under reflection and the four
coordinate quarter-turns. But
\[
 B(1,1)=1,\qquad 1^2+1^2=2\ne B(1,1)^2.
\]
So a finite collection of elementary symmetries cannot
by itself force the Euclidean length formula.

Lean proves more: among matrices of the form~\eqref{eq:normal},
the only $B$-rigid ones are the four signed coordinate
quarter-turns. A real additive action containing $Q$
would also contain a rigid square root $F_{\theta/2}$
of $Q$; no such matrix exists for $B$. Hence no
$B$-rigid additive real linear action containing
the primitive quarter-turn exists, \emph{even without}
a continuity hypothesis. The formal statements include
\texttt{boxLength\_rigid\_normal\_iff},
\texttt{boxLength\_not\_pythagorean}, and
\texttt{boxLength\_no\_additive\_quarter\_turn\_action}.
This example is an obstruction to \emph{finite}
symmetries, not a counterexample satisfying all
the hypotheses of either main theorem.

\section{Infinitesimal interpretation and proof boundaries}

If the motion coefficients have the stated derivatives
at zero, the infinitesimal generator has the form
$\omega J$ after the radial derivative is shown to
vanish. Algebraically, for a symmetric matrix
$M$, the infinitesimal invariance equation
$J^TM+MJ=0$ forces $M=\lambda I$.
The repository contains separate Lean statements on
the generator, conserved quadratic expressions, and
the classification of infinitesimally invariant
symmetric forms. These \emph{supplement} the group
argument above; differentiability and Lie theory are
not needed for Theorem~\ref{thm:main}.

The precise boundary matters. Lean establishes
that the displayed conclusions follow from formal
hypotheses, and the explicit model shows joint
satisfiability in standard real analysis. It does
not prove that the action exists from bare incidence
or primitive congruence axioms, establish a minimal
axiom system, or certify historical novelty. The
principle that rich linear isometry groups select
Euclidean structure has substantial classical context.
A kernel-level axiom audit concerns Lean's logical
dependencies, not whether these \emph{geometric}
hypotheses are minimal.

\section{Where to find the formal proofs}

The proof steps above correspond to the following
files in the \texttt{palomar-preparation} branch:
\begin{itemize}
 \item \texttt{PythagoreanIsotropy.lean}: primitive quarter-turn,
       motion normal form, and the archived conditional theorem;
 \item \texttt{ContinuousIsotropy.lean}: derived axis
       reachability;
 \item \path{CompactCharacter.lean}: the periodic-character
       argument;
 \item \path{CompactIsotropy.lean}: quadratic preservation
       and length identification without mirror symmetry;
 \item \texttt{EuclideanModel.lean} and
       \path{IsotropyCharacterization.lean}: explicit model
       and the characterization equivalence;
 \item \path{IsotropyCounterexamples.lean}: maximum-norm
       obstruction.
\end{itemize}
\noindent\textbf{Repository:}
\href{https://github.com/alexodas/pythagorean-isotropy-lean/tree/palomar-preparation}{pythagorean-isotropy-lean (palomar-preparation branch)}.
The root \texttt{README.md} provides reproducible toolchain,
Comparator, verification, archival and possible Palomar
publication details; those operational details are deliberately
kept out of this mathematical exposition.

\end{document}
